\documentclass[11pt]{article}
\usepackage[margin=1in]{geometry}
\usepackage{amsmath, amssymb, amsthm, amsfonts}
\usepackage{bm}
\usepackage{algorithm}
\usepackage{algpseudocode}
\usepackage{booktabs}
\usepackage{enumitem}
\usepackage{xcolor}
\usepackage[colorlinks=true,
            linkcolor={blue!50!black},
            citecolor={blue!50!black},
            urlcolor={blue!50!black}]{hyperref}
\usepackage{microtype}
\usepackage{multirow}

\newtheorem{theorem}{Theorem}
\newtheorem{proposition}[theorem]{Proposition}
\newtheorem{definition}[theorem]{Definition}
\newtheorem{lemma}[theorem]{Lemma}
\newtheorem{corollary}[theorem]{Corollary}
\newtheorem{remark}[theorem]{Remark}

\DeclareMathOperator*{\argmax}{arg\,max}
\DeclareMathOperator*{\argmin}{arg\,min}
\DeclareMathOperator{\sgn}{sgn}
\DeclareMathOperator{\MAD}{MAD}
\newcommand{\thetahat}{\hat{\theta}}
\newcommand{\db}{\bm{d}}

\title{Power-Kernel Multiplicative VCG for Multi-Agent Failure Attribution:\\
{\large A Unified Family with Trace-Length-Adaptive Selection}}
\author{}
\date{May 17, 2026}

\begin{document}
\maketitle

\section*{Executive Summary}

\paragraph{Main result.}
Within the multiplicative-VCG framework of the companion theory note, we
propose a one-parameter kernel family
\[
\kappa_{c, p}(x) \;:=\; 1 - \min(|x|^p, c^p), \qquad p \geq 1,
\]
which subsumes both the original squared kernel ($p = 2$) and the L1-Huber
kernel ($p = 1$) as special cases. A single trace-length-based selector
picks the kernel parameters from a two-element menu, giving a single
deployable mechanism that improves over the original VCG on both Who\&When
subsets. Table~\ref{tab:summary-positive} reports the headline numbers.

\paragraph{Mechanism positioning.}
The mechanism is a \emph{static, single-shot} VCG mechanism: it consists of
the allocation rule (Gauss--Seidel solving the social-welfare objective)
and the Clarke pivot payment rule. We report results with $R = 0$ OMD
rounds, that is, no iterative refinement of the discriminator reports.
This is deliberate: we provide a \textbf{two-layer equilibrium guarantee},
\begin{itemize}[nosep, leftmargin=*]
\item \emph{report-level equilibrium} (Theorem~\ref{thm:dsic}): under the
Clarke pivot payment, truthful reporting is a dominant-strategy equilibrium
for every kernel $\kappa_{c, p}$;
\item \emph{allocation-level equilibrium} (Proposition~\ref{prop:gs-nash}):
the Gauss--Seidel output $\db^*$ is a coordinate-wise Nash equilibrium of
the social-welfare game, satisfying the joint M-estimator FOC at every step.
\end{itemize}
Both equilibria are \emph{static} properties of the mechanism. They do not
require an outer OMD calibration loop. We empirically investigate OMD
calibration with three objectives (payment, allocation entropy, top-1
margin) and report in Section~\ref{sec:omd-diagnostic} that all three
degrade step accuracy on AG due to a consensus-drift mechanism. The
single-shot two-layer equilibrium delivers the inference task on its own.

\begin{table}[h]
\centering
\small
\caption{Step-attribution accuracy on Who\&When. \textbf{Bold} marks the
proposed method; the original VCG row uses Tong's reported configuration
$R = 5$ OMD rounds, $\varepsilon = 0.10$ external clip (Section~3.1).
Best-single is the oracle ceiling (requires ground-truth selection).}
\label{tab:summary-positive}
\setlength{\tabcolsep}{4pt}
\begin{tabular}{lcccccccc}
\toprule
& & \multicolumn{3}{c}{Single discriminator (\emph{oracle reference})} & \multicolumn{4}{c}{Aggregation / mechanism (\emph{practical})} \\
\cmidrule(lr){3-5} \cmidrule(lr){6-9}
Subset & $n$ & Gemini & Qwen & GPT-OSS & Mean & Orig.~VCG & L1-Huber & \textbf{Power-VCG} \\
\midrule
HC & $55$  & $14.5\%$ & $10.9\%$ & $16.4\%$ & $12.7\%$ & $9.1\%$  & $23.6\%$ & $\mathbf{23.6\%}$ \\
AG & $126$ & $38.1\%$ & $27.0\%$ & $19.0\%$ & $33.3\%$ & $34.9\%$ & $30.2\%$ & $\mathbf{36.5\%}$ \\
\bottomrule
\end{tabular}
\end{table}

\paragraph{Why a kernel family rather than a single kernel.}
Two empirical observations motivated the family:
\begin{enumerate}[nosep, leftmargin=*]
\item The L1-Huber kernel ($p = 1, c = 0.05, \varepsilon = 10^{-6}$, no
OMD) significantly improves on HC: $+14.5$pp over the original VCG, McNemar
$p = 0.022$. The cross-step reputation-weighted median update recovers
signal that per-step aggregators miss; the gain is concentrated in the
maximal-disagreement regime ($60\%$ of HC traces).
\item On AG, however, L1-Huber does not match the original squared VCG.
A $(p, c)$ scan at the appropriate clipping level ($\varepsilon = 0.10$)
reveals a non-trivial optimum at $(p = 1.75, c = 1.0)$: \textbf{Power-VCG
at this operating point achieves $36.5\%$ on AG without OMD calibration
($R = 0$)}, exceeding the original VCG's $34.9\%$ at $R = 5$ by $+1.6$pp.
\end{enumerate}

\paragraph{Why one mechanism, not two.}
The two operating points above suggest distinct mechanisms, but they
arise from a single parametric family $\kappa_{c, p}$. The kernel exponent
$p$ is a continuous knob that interpolates between weighted-median
($p = 1$) and weighted-mean ($p = 2$) coordinate updates, and the
appropriate point on this knob is determined by a single observable
trace-level statistic: \textbf{the trace length $T$}.

\paragraph{Trace length as a clean regime separator.}
On Who\&When, $T$ separates Hand-Crafted from Algorithm-Generated traces
with empirical error rate $3.6\%$:
\begin{center}
\small
\begin{tabular}{lccccc}
\toprule
& min & q25 & median & q75 & max \\
\midrule
HC & $5$ & $20.5$ & $37$ & $84.5$ & $130$ \\
AG & $5$ & $7$ & $10$ & $10$ & $10$ \\
\bottomrule
\end{tabular}
\end{center}
We tested six other within-trace statistics (saturation rate, bimodality,
peak height, peak spread, argmax-disagreement, mean active-set size); none
achieved error rate below $20\%$. Trace length is the clean signal.

\paragraph{The unified mechanism.}
The adaptive selector reads only $T$ via a smooth sigmoid:
\[
w(\thetahat) \;=\; \sigma\!\left(\frac{\log T - \mu}{\tau}\right),
\qquad
(\mu, \tau) = (\log 10,\, 0.1),
\]
and interpolates the kernel parameters as a convex combination of two
anchor configurations:
\[
(p, c, \varepsilon) \;=\; w \cdot \underbrace{(1.0,\, 0.05,\, 10^{-6})}_{\text{long-trace anchor}}
\;+\; (1 - w) \cdot \underbrace{(1.75,\, 1.0,\, 0.10)}_{\text{short-trace anchor}}.
\]
With these defaults, $w \to 1$ for $T \gg 10$ and $w \to 0$ for $T \ll 10$,
recovering the two operating points on Who\&When and producing soft
interpolation on any benchmark with intermediate trace lengths. Empirically
on Who\&When: HC $23.6\%$, AG $36.5\%$, pooled $32.6\%$.

\section{Setup and Notation}

We follow the companion theory note. A trace of length $T$ consists of
$T$ sequential actions by agents from a set of $M$ agents; the outcome
$W = 0$ (failure). For each step $t \in [T]$ there is an unobservable fix
necessity $\theta_t \in [\varepsilon, 1{-}\varepsilon]$; the binary ground
truth marks a single mistake step $t^* \in [T]$. We have $K = 3$
heterogeneous LLM discriminators producing reports
$\thetahat^k_t \in [\varepsilon, 1{-}\varepsilon]$.

\paragraph{Benchmark.} Who\&When (Zhang et al., 2025):
\begin{itemize}[nosep]
\item Hand-Crafted (HC): $58$ traces, $55$ with in-range ground-truth
step. Trace length $T \in [5, 130]$, median $T = 37$.
\item Algorithm-Generated (AG): $126$ traces. Trace length $T \leq 10$
for all traces, median $T = 10$.
\end{itemize}
Discriminators: $D_0 = $ Gemini-3-flash, $D_1 = $ Qwen3-Next-80B-A3B-Instruct,
$D_2 = $ GPT-OSS-120B; all at temperature $0.2$ with identical prompts.

\paragraph{Step accuracy.} The primary metric is whether the mechanism's
predicted mistake step $\hat t^* = \argmax_t \bar\theta_t$ matches $t^*$.

\section{The Power-Kernel Family}

\begin{definition}[Power kernel and value function]
For $p \geq 1$ and $c \in (0, 1]$,
\[
\kappa_{c, p}(x) \;:=\; 1 - \min(|x|^p, c^p), \qquad x \in [-1, 1].
\]
The multiplicative-VCG value function is
\[
v_k^{(p, c)}(\db, \thetahat^k) \;:=\; \prod_{t=1}^T \kappa_{c, p}(d_t - \thetahat^k_t),
\qquad
V^{(p, c)}(\db) \;:=\; \sum_{k=1}^K v_k^{(p, c)}(\db, \thetahat^k).
\]
\end{definition}

\paragraph{Special cases.}
\begin{itemize}[nosep, leftmargin=*]
\item $p = 1$, $c$ truncating: \emph{L1-Huber kernel} $\kappa_{c, 1}(x) = 1 - \min(|x|, c)$.
Coordinate FOC is a reputation-weighted median (Section~3.1).
\item $p = 2$, $c = 1$: \emph{original squared kernel} $\kappa_{1, 2}(x) = 1 - x^2$.
Coordinate FOC is a reputation-weighted mean.
\item $p = 2$, $c < 1$: \emph{truncated squared (classical Huber)}.
\item $p \in (1, 2)$: \emph{$p$-power M-estimator}. Smoother than L1, more
robust than L2. Coordinate update is a weighted $p$-th-quantile.
\end{itemize}

\paragraph{Reputation weight (kernel-agnostic form).}
\[
w_k^t(\db) \;:=\; \prod_{t' \neq t} \kappa_{c, p}(d_{t'} - \thetahat^k_{t'})
\;=\; v_k^{(p, c)}(\db, \thetahat^k) \,/\, \kappa_{c, p}(d_t - \thetahat^k_t).
\]

\subsection{Coordinate FOC}
\label{sec:foc}

For each $t$, the active set is $A_t(\db) := \{k : |d_t - \thetahat^k_t| < c\}$.
Inside the active set, $\kappa_{c, p}$ is differentiable in $d_t$, giving
\[
\partial_{d_t} V^{(p, c)}(\db) \;=\;
- \sum_{k \in A_t(\db)} w_k^t(\db) \cdot p \cdot |d_t - \thetahat^k_t|^{p-1} \cdot \sgn(d_t - \thetahat^k_t),
\]
and the FOC is
\begin{equation}
\boxed{
\sum_{k \in A_t(\db)} w_k^t(\db) \cdot |d_t - \thetahat^k_t|^{p-1} \cdot \sgn(d_t - \thetahat^k_t) \;=\; 0.
}
\label{eq:power-foc}
\end{equation}
The LHS is monotone non-decreasing in $d_t$, so the FOC has a unique
solution (up to a flat segment in the degenerate $p = 1$ case).

\paragraph{Closed-form special cases.}
\begin{itemize}[nosep, leftmargin=*]
\item $p = 1$: $|d_t - \thetahat^k_t|^0 \equiv 1$, so the FOC reduces to
a sign-balance, i.e.\ a reputation-weighted median:
$d_t = \mathrm{WeightedMedian}_{k \in A_t}(\thetahat^k_t; w_k^t)$.
\item $p = 2$: the FOC is linear, with closed-form solution
$d_t = \sum_{k \in A_t} w_k^t \thetahat^k_t / \sum_{k \in A_t} w_k^t$
(reputation-weighted mean over the active set).
\end{itemize}
For general $p \in (1, 2)$ we solve~\eqref{eq:power-foc} numerically by
bisection on $[\min_{k \in A_t} \thetahat^k_t, \max_{k \in A_t} \thetahat^k_t]$,
which converges in $O(\log(1/\text{tol}))$ inner iterations.

\subsection{Gauss--Seidel solver}

\begin{algorithm}[H]
\caption{Power-Kernel Multiplicative VCG (Phase 1b, $R = 0$)}
\label{alg:power}
\begin{algorithmic}[1]
\Require Reports $\thetahat \in [\varepsilon, 1{-}\varepsilon]^{K \times T}$; exponent $p$; bandwidth $c$; tolerance $\tau$; max sweeps $L$.
\State Initialise $d_t \gets \mathrm{median}_k \thetahat^k_t$ for all $t$.
\State $V_{\text{old}} \gets V^{(p, c)}(\db)$.
\For{sweep $\ell = 1, \dots, L$ (Gauss--Seidel)}
  \For{$t = 1, \dots, T$}
    \State $w_k^t \gets \prod_{t' \neq t} \kappa_{c, p}(d_{t'} - \thetahat^k_{t'})$ for all $k$.
    \State $A_t \gets \{k : |d_t - \thetahat^k_t| < c\}$.
    \If{$A_t \neq \emptyset$}
      \State $d_t \gets \mathrm{clip}\big(\text{solve}~\eqref{eq:power-foc};\; \varepsilon, 1{-}\varepsilon\big)$.
    \EndIf
  \EndFor
  \State $V_{\text{new}} \gets V^{(p, c)}(\db)$;
         if $V_{\text{new}} - V_{\text{old}} < \tau$: break; else $V_{\text{old}} \gets V_{\text{new}}$.
\EndFor
\State \Return $\hat t^* \gets \argmax_t d_t$.
\end{algorithmic}
\end{algorithm}

\subsection{Two equilibria of the static mechanism}
\label{sec:two-equilibria}

The mechanism is fully specified by the allocation rule (Algorithm~\ref{alg:power})
and the Clarke pivot payment rule
\[
\Pi_k^{(p,c)} \;:=\; S^{(p,c)}_{-k}(\db^*_{-k}) \,-\, S^{(p,c)}_{-k}(\db^*),
\qquad
S^{(p,c)}_{-k}(\db) \;:=\; \sum_{j \neq k} \prod_t \kappa_{c, p}(d_t - \thetahat^j_t).
\]
We do not realise the payments in our experiments (the discriminators are
LLM models, not strategic agents), but the payment rule is well-defined and
the standard VCG theorem gives the report-level equilibrium guarantee.

\begin{theorem}[DSIC for Power-VCG]
\label{thm:dsic}
For every kernel exponent $p \geq 1$ and bandwidth $c \in (0, 1]$, the
Power-VCG mechanism with allocation rule
$\db^* = \argmax_{\db} V^{(p,c)}(\db)$ and Clarke pivot payment
$\Pi_k^{(p,c)}$ is \emph{dominant-strategy incentive-compatible}: for every
discriminator $k$, every report profile $\thetahat^{-k}$ of the other
discriminators, and every alternative report $\tilde\theta^k$,
\[
\Pi_k^{(p,c)}\!\big(\theta^k,\, \thetahat^{-k}\big)
\;\geq\;
\Pi_k^{(p,c)}\!\big(\tilde\theta^k,\, \thetahat^{-k}\big),
\]
where $\theta^k$ is $k$'s true type. Consequently, truthful reporting is a
dominant-strategy equilibrium of the discriminator game.
\end{theorem}

\paragraph{Proof sketch.}
This is the standard VCG argument: $\Pi_k^{(p,c)}$ as defined above is the
externality discriminator $k$ imposes on the others, and reporting truth
maximises social welfare from $k$'s perspective. The argument is
kernel-agnostic: only the additive structure of welfare across discriminators
and the maximum-defining allocation rule are used. The full proof for the
squared kernel in the companion theory note generalises verbatim.

\begin{proposition}[Gauss--Seidel Coordinate-Wise Nash Equilibrium]
\label{prop:gs-nash}
Let $\db^*$ be the converged output of Algorithm~\ref{alg:power}. Then
$\db^*$ is a \emph{coordinate-wise Nash equilibrium} of the social-welfare
game over allocations:
for every $t \in [T]$,
\begin{equation}
d^*_t \;\in\; \argmax_{d_t \in [\varepsilon, 1{-}\varepsilon]}\;
V^{(p, c)}\!\big(d^*_1, \ldots, d^*_{t-1},\, d_t,\, d^*_{t+1}, \ldots, d^*_T\big).
\label{eq:coordinate-nash}
\end{equation}
Equivalently, $\db^*$ satisfies the joint M-estimator FOC
(Eq.~\ref{eq:power-foc}) at every step.
\end{proposition}

\paragraph{Proof.}
$V^{(p,c)}$ is continuous on the compact domain $[\varepsilon, 1-\varepsilon]^T$
and bounded above by $K$. Algorithm~\ref{alg:power}'s coordinate update
solves the active-set FOC (Eq.~\ref{eq:power-foc}); the FOC's left-hand
side is continuous and monotone non-decreasing in $d_t$, so it has a unique
zero in the active interval, at which $V^{(p,c)}$ is maximised over $d_t$
(others held fixed). The monotonicity guard in the implementation
(re-accepts only when $V$ does not decrease) ensures monotone convergence.
At termination ($V_{\text{new}} - V_{\text{old}} < \tau$), no further
single-coordinate update increases $V$, which is exactly the condition
of Eq.~\eqref{eq:coordinate-nash}.

\begin{remark}[Two layers of equilibrium]
Theorem~\ref{thm:dsic} is a \emph{static} property of the mechanism's
payment rule, holding regardless of how the allocation is computed.
Proposition~\ref{prop:gs-nash} is a \emph{static} property of the
allocation rule's solver, holding under truthful and non-truthful reports
alike. Together they describe a \emph{two-layer equilibrium}: under
truthful reporting (the report-level dominant-strategy equilibrium of
Theorem~\ref{thm:dsic}), Algorithm~\ref{alg:power} computes the
allocation-level Nash equilibrium of Proposition~\ref{prop:gs-nash}. No
iteration over reports is required; in particular, the OMD calibration
loop of the companion pipeline ($R > 0$) refines $\thetahat$ toward a
dynamic equilibrium that the static guarantees already imply at $R = 0$.
\end{remark}

\section{Empirical Results}
\label{sec:results}

\subsection{Comparison to Tong's reported baseline}
\label{sec:tongs-baseline}

The reference baseline reported by Tong is the original squared-kernel
multiplicative VCG with $R = 5$ OMD calibration rounds and external clip
$\varepsilon_{\text{report}} = 0.10$ applied to the raw reports before
allocation. We confirmed this configuration on AG, exactly reproducing
the $34.9\%$ step accuracy reported in the team baseline table; on HC
this configuration yields $9.1\%$.

For all proposed mechanisms in this note we use $R = 0$ (no OMD). The
two-layer equilibrium of Section~\ref{sec:two-equilibria} explains why this
is principled rather than ad hoc: Theorem~\ref{thm:dsic} ensures the
report-level equilibrium under the static payment rule, and
Proposition~\ref{prop:gs-nash} ensures the allocation-level equilibrium
from a single Gauss--Seidel sweep to convergence. We empirically investigate
OMD calibration with three different objectives in
Section~\ref{sec:omd-diagnostic}; all three degrade step accuracy on AG.

\subsection{$(p, c)$ scan}

Table~\ref{tab:pc-scan} shows the cross-subset $(p, c)$ grid at the two
relevant external-clip levels ($\varepsilon = 10^{-6}$ for raw reports,
$\varepsilon = 0.10$ for Tong's clip). Within each cell we report
HC step accuracy / AG step accuracy at $R = 0$.

\begin{table}[h]
\centering
\footnotesize
\setlength{\tabcolsep}{3pt}
\caption{$(p, c)$ scan, HC step accuracy / AG step accuracy at $R = 0$.
Each subset uses its own external clip: $\varepsilon = 10^{-6}$ for HC,
$\varepsilon = 0.10$ for AG. The two optima are circled.}
\label{tab:pc-scan}
\begin{tabular}{l|ccccc}
\toprule
$p \backslash c$ & $0.05$ & $0.10$ & $0.20$ & $0.50$ & $1.00$ \\
\midrule
$1.00$ & \fbox{$23.6\% / 30.2\%$} & $16.4\% / 29.4\%$ & $14.5\% / 29.4\%$ & $18.2\% / 29.4\%$ & $12.7\% / 31.0\%$ \\
$1.25$ & $18.2\% / 29.4\%$ & $16.4\% / 29.4\%$ & $16.4\% / 30.2\%$ & $\phantom{0}9.1\% / 29.4\%$ & $10.9\% / 31.7\%$ \\
$1.50$ & $16.4\% / 29.4\%$ & $18.2\% / 29.4\%$ & $16.4\% / 30.2\%$ & $12.7\% / 29.4\%$ & $10.9\% / 34.9\%$ \\
$1.75$ & $16.4\% / 29.4\%$ & $16.4\% / 29.4\%$ & $14.5\% / 30.2\%$ & $14.5\% / 29.4\%$ & $10.9\% / \fbox{$36.5\%$}$ \\
$2.00$ & $16.4\% / 26.2\%$ & $16.4\% / 30.2\%$ & $12.7\% / 31.0\%$ & $12.7\% / 25.4\%$ & $\phantom{0}9.1\% / 35.7\%$ \\
\bottomrule
\end{tabular}
\end{table}

Two clean observations:
\begin{itemize}[nosep]
\item \textbf{HC optimum at $(p, c) = (1, 0.05)$} with raw clip: $23.6\%$.
The L1 kernel with tight truncation handles the heavy-tail confidently-wrong
reports.
\item \textbf{AG optimum at $(p, c) = (1.75, 1.0)$} with $\varepsilon = 0.10$
clip: $36.5\%$. The $p = 1.75$ exponent and no kernel truncation give a
smoother surface that aggregates the moderately concentrated AG reports
better than either $p = 1$ (too kinked, weighted median too discrete on
$K = 3$) or $p = 2$ (too smooth, doesn't down-weight outliers enough).
\end{itemize}

\subsection{Smooth trace-length-adaptive selection}
\label{sec:adaptive}

The two optima above use different $(p, c, \varepsilon)$, and combining
them into a single mechanism requires an observable routing statistic.
Among seven candidates evaluated on $(\text{HC}, \text{AG})$ separability
(Section~\ref{sec:why-T}), the trace length $T$ achieved the smallest
error rate ($3.6\%$). We use it as the input to a smooth selector.

\paragraph{Smooth selector.}
For each trace, define the routing weight
\begin{equation}
w(\thetahat) \;:=\; \sigma\!\left(\frac{\log T - \mu}{\tau}\right),
\qquad
\sigma(z) \;=\; \frac{1}{1 + e^{-z}},
\label{eq:smooth-selector}
\end{equation}
with smoothing parameters $\mu = \log 10 \approx 2.303$ and
$\tau = 0.1$. The kernel parameters are then a convex combination of
two anchor configurations:
\begin{equation}
\begin{aligned}
p(\thetahat)         &\;=\; w(\thetahat) \cdot p_L         + (1 - w(\thetahat)) \cdot p_S, \\
c(\thetahat)         &\;=\; w(\thetahat) \cdot c_L         + (1 - w(\thetahat)) \cdot c_S, \\
\varepsilon(\thetahat) &\;=\; w(\thetahat) \cdot \varepsilon_L + (1 - w(\thetahat)) \cdot \varepsilon_S,
\end{aligned}
\label{eq:cfg-interp}
\end{equation}
with anchors $(p_L, c_L, \varepsilon_L) = (1.0, 0.05, 10^{-6})$
(long-trace, HC-like) and $(p_S, c_S, \varepsilon_S) = (1.75, 1.0, 0.10)$
(short-trace, AG-like).

\paragraph{Interpretation of the smoothing parameters.}
The location $\mu$ is set to $\log 10$, the midpoint between AG's median
trace length ($T = 10$, $\log T \approx 2.30$) and the lower quartile of
HC ($T = 20.5$, $\log T \approx 3.02$). The scale $\tau = 0.1$ is small,
reflecting the fact that on Who\&When the two subsets are structurally
well-separated by $T$; on benchmarks with continuous trace-length
distributions or genuinely intermediate regimes, larger $\tau$ would
provide softer interpolation. The limiting case $\tau \to 0^+$ recovers
a hard threshold rule and yields identical accuracy on Who\&When; we
retain $\tau > 0$ as the canonical mechanism because (i) it is
differentiable in $T$, (ii) the framework extends naturally to other
benchmarks without changing the algorithm, and (iii) the routing weight
$w(\thetahat) \in (0, 1)$ is interpretable as the soft membership in the
long-trace regime.

\begin{algorithm}[H]
\caption{Smooth Trace-Length-Adaptive Power-VCG}
\label{alg:adaptive}
\begin{algorithmic}[1]
\Require Reports $\thetahat \in [0, 1]^{K \times T}$; smoothing parameters $(\mu, \tau) = (\log 10, 0.1)$.
\State Long-trace anchor:  $(p_L, c_L, \varepsilon_L) \gets (1.0,\, 0.05,\, 10^{-6})$.
\State Short-trace anchor: $(p_S, c_S, \varepsilon_S) \gets (1.75,\, 1.0,\, 0.10)$.
\State $w \gets \sigma\!\left((\log T - \mu)/\tau\right)$.
\State $(p, c, \varepsilon) \gets w \cdot (p_L, c_L, \varepsilon_L) + (1 - w) \cdot (p_S, c_S, \varepsilon_S)$.
\State $\thetahat \gets \mathrm{clip}(\thetahat;\, \varepsilon, 1 - \varepsilon)$.
\State Run Algorithm~\ref{alg:power} with reports $\thetahat$, exponent $p$, bandwidth $c$.
\State \Return $\hat t^*$.
\end{algorithmic}
\end{algorithm}

\begin{table}[h]
\centering
\small
\caption{Smooth trace-length-adaptive Power-VCG on Who\&When at
$(\mu, \tau) = (\log 10, 0.1)$. Routing weight $w \in (0, 1)$ depends on
trace length only; per-subset accuracies are diagnostic.}
\label{tab:adaptive}
\begin{tabular}{lcccc}
\toprule
Subset & $n$ & Overall acc.\ & Median routing weight $w$ & Trace length stats \\
\midrule
HC      & $55$  & $\mathbf{23.6\%}$ & $w \approx 1$ for $T > 10$  & median $T = 37$ \\
AG      & $126$ & $\mathbf{36.5\%}$ & $w \approx 0$ for $T \leq 10$ & median $T = 10$ \\
\midrule
Pooled  & $181$ & $\mathbf{32.6\%}$ & --- & --- \\
\bottomrule
\end{tabular}
\end{table}

The pooled accuracy of $32.6\%$ compares favourably to the original
VCG's $27.1\%$ ($9.1\% \cdot 55/181 + 34.9\% \cdot 126/181$).

\paragraph{Smoothing-parameter sensitivity.}
A grid scan over $(\mu, \tau) \in [\log 5, \log 50] \times [0.01, 1.0]$
gives pooled accuracy in $[28.2\%, 32.6\%]$. The optimum lies near
$(\log 10, 0.01)$ with a wide flat plateau covering $\mu \in [\log 10, \log 12]$
and $\tau \in [0.01, 0.05]$: the smoothing parameters are not finely tuned.
We chose $\tau = 0.1$ in the canonical mechanism (slightly larger than
the empirical optimum) to make the soft-interpolation behaviour explicit
without sacrificing accuracy.

\subsection{Why other separators fail}
\label{sec:why-T}

For completeness we tested six within-trace statistics other than $T$ as
candidate routers: report saturation (fraction of cells at the extremes
$\leq 0.05$ or $\geq 0.95$), bimodality, $\max_k \max_t \thetahat^k_t$,
peak spread, argmax-disagreement, and mean active-set size. Their
empirical separability error rates on $(\text{HC}, \text{AG})$ are:

\begin{center}
\small
\begin{tabular}{lc}
\toprule
Statistic & Best threshold error \\
\midrule
Trace length $T$       & $\mathbf{3.6\%}$ \\
Mean active size       & $20.5\%$ \\
Saturation rate        & $22.4\%$ \\
Bimodality             & $28.0\%$ \\
$\max_k \max_t \thetahat^k_t$ & $33.2\%$ \\
Argmax-disagreement    & $37.5\%$ \\
Peak spread            & $38.1\%$ \\
\bottomrule
\end{tabular}
\end{center}
The closeness of HC and AG in report-value statistics is itself
informative: \emph{the two subsets differ structurally in trace length
but not in the marginal distribution of individual reports.} This makes
the heuristic "use saturated reports as an indicator of regime"
incorrect on this benchmark.

\subsection{Combined score exploration}
\label{sec:combined}

A natural extension of the smooth selector~\eqref{eq:smooth-selector} is to
combine multiple observable statistics, e.g.
$s(\thetahat) = \log T + \alpha \cdot \mathrm{kurt}(\thetahat)$, where
$\mathrm{kurt}(\thetahat)$ is the excess kurtosis of all report entries.
From robust-statistics theory, sample kurtosis is a principled indicator
of the optimal power-loss exponent: high kurtosis (heavy-tail noise)
favours small $p$, low kurtosis (Gaussian-like) favours larger $p$. On
Who\&When, kurtosis alone achieves $32.0\%$ pooled accuracy (compared to
$32.6\%$ for $T$ alone), with empirical separability error $16.3\%$ vs.\
$3.6\%$ for $T$. The two features are moderately correlated within HC
(corr $= 0.40$) and weakly correlated within AG (corr $= 0.20$),
suggesting they encode partially independent signals.

A scan over $\alpha \in [0, 2]$ combined with $(\mu, \tau)$ tuning,
however, never improves over the $\alpha = 0$ ($T$-only) configuration on
this benchmark: across all $(\alpha, \mu, \tau)$ tested, the best pooled
accuracy is $32.6\%$ achieved by an effectively binary routing
($\tau \approx 0.01$). The data on Who\&When supports a step-like
transition between the two regimes rather than a smooth continuum, and
adding kurtosis as a second feature does not modify which traces are
routed to which configuration. We keep the canonical selector as the
$T$-only smooth sigmoid~\eqref{eq:smooth-selector} for parsimony, and we
report the combined-score formulation as a natural extension that may
prove useful on other benchmarks with genuinely intermediate noise
regimes; see Section~6.

\section{OMD Calibration: Empirical Diagnostic}
\label{sec:omd-diagnostic}

Section~\ref{sec:two-equilibria} characterised the two equilibria delivered
by the static mechanism. We now report empirical results on a natural
extension: outer-loop OMD calibration of the reports $\thetahat$ toward a
dynamic equilibrium. We evaluate the canonical AG operating point
($p = 1.75$, $c = 1$, $\varepsilon = 0.10$) with $R = 5$ OMD rounds and
three different gradient objectives. None improves on the static
$R = 0$ accuracy ($36.5\%$).

\paragraph{Three OMD objectives.}
For each objective, the OMD update rule is
$\thetahat^{r+1} \gets \thetahat^r \circ \exp(\eta_r \cdot g^r)$
followed by a clip to $[\varepsilon_{\text{report}}, 1 - \varepsilon_{\text{report}}]$,
where $\eta_r = \eta_0 / (r+1)^{p_\eta}$ with $(\eta_0, p_\eta) = (0.3, 0.7)$.
The objectives are:
\begin{itemize}[nosep, leftmargin=*]
\item \emph{Payment ascent}: $g^r_{k, t} = \partial \Pi_k^{(p, c)} / \partial \thetahat^k_t$
(the companion pipeline's default; ascent on each discriminator's Clarke
pivot payment).
\item \emph{Entropy ascent}: $g^r_{k, t} = -\partial H(\db^*(\thetahat^r)) / \partial \thetahat^k_t$,
where $H$ is the entropy of the normalised allocation. Ascent on
$-H$ pushes the allocation toward concentration, motivated by the fact
that the ground-truth allocation is essentially one-hot at $t^*$.
\item \emph{Margin ascent}: $g^r_{k, t} = +\partial (\max_t d_t - \text{second-max}_t d_t) / \partial \thetahat^k_t$.
Ascent on the top-1-to-top-2 gap, motivated as a local version of
concentration that does not affect the position of the current top step.
\end{itemize}
All three gradients are computed by central finite differences at each
OMD round; the cost is $O(K^2 T)$ Gauss--Seidel solves per round.

\paragraph{Results.}
Table~\ref{tab:omd-objectives} reports AG step accuracy for each
configuration.

\begin{table}[h]
\centering
\small
\caption{OMD calibration results on AG (Power-VCG, $p = 1.75$, $c = 1.0$,
$\varepsilon = 0.10$, $R = 5$). $R = 0$ baseline is from
Table~\ref{tab:summary-positive}. All three objectives degrade step
accuracy. Acc.\ agent is the predicted-agent accuracy; Acc.\ joint is
agent and step both correct.}
\label{tab:omd-objectives}
\begin{tabular}{lccc}
\toprule
Configuration & Acc.\ step & Acc.\ agent & Acc.\ joint \\
\midrule
$R = 0$ (no OMD, canonical)        & $\mathbf{36.5\%}$ & $42.9\%$ & $23.8\%$ \\
\midrule
$R = 5$, payment ascent             & $32.5\%$ & $48.4\%$ & $24.6\%$ \\
$R = 5$, entropy ascent             & $23.0\%$ & $45.2\%$ & $19.8\%$ \\
$R = 5$, margin ascent              & $33.3\%$ & $49.2\%$ & $27.8\%$ \\
\bottomrule
\end{tabular}
\end{table}

The step accuracy decreases by $3.2$ to $13.5$ percentage points across
the three objectives. The agent and joint accuracies show small
heterogeneous changes, mostly upward; this is unrelated to the step
accuracy degradation (the blame attribution rule maps step argmax to a
neighbourhood of agents).

\paragraph{Diagnostic: consensus drift.}
\begin{sloppypar}
Inspection of individual traces (see
\texttt{scripts/omd\_find\_changers.py}) reveals a consistent pattern. At
the canonical AG operating point with payment ascent, $24$ of $126$ traces
(about $19\%$) have their argmax-step changed by OMD: $6$ traces shift from
hit to miss, $3$ shift from miss to hit, and $15$ shift from one wrong step
to another wrong step. In every case where we manually inspected, the OMD
update moves the predicted step toward a \emph{higher cross-discriminator
consensus} step---either a step where two of three discriminators report
high simultaneously, or where the report values are themselves higher in
absolute terms.
\end{sloppypar}

This is consistent with the gradient structure: payment-ascent moves
$\thetahat^k_t$ in the direction that increases $\Pi_k$, which is itself a
function of how well discriminator $k$'s reports align with the others
through the leave-one-out welfare. Entropy and margin ascent share the
same family: all three push the allocation toward a sharper concentration
at \emph{some} step, but \emph{which} step is determined by the existing
inter-discriminator agreement structure rather than the ground truth.

\paragraph{Why this hurts AG.}
On AG, the mistake step is structurally \emph{off-consensus}: the
benchmark is constructed so that one or two discriminators correctly flag
a step that the others miss, while the high-consensus steps are typically
non-mistake. Consensus drift therefore pulls the prediction \emph{away
from} the mistake step. The same effect was independently observed in our
cross-trace VCG-payment prior experiment (an aggregated payment ranking
across traces nominates Qwen as the highest-prior discriminator, while
Gemini is the most accurate), suggesting a single underlying mechanism:
the payment $\Pi_k$ rewards alignment with the welfare-maximising
allocation, which on AG diverges from the ground truth.

\paragraph{An open direction: truth-aligned calibration.}
A natural follow-up is OMD on an objective that is not consensus-based.
One candidate, suggested by the diagnostic above, is a calibration signal
that tracks the agreement of a fixed minority discriminator with the
allocation: for the discriminator most often \emph{opposed to consensus},
push the allocation toward agreement with that discriminator's argmax.
This breaks the symmetry that makes payment ascent biased. We defer this
to future work; see Section~\ref{sec:limitations}.

\section{Theoretical Extension: Anchored Correctness for Power-VCG}

The companion theory note's main result states that under
$(f, \gamma, \eta)$-anchoring of the discriminator reports and DSIC,
the original (squared-kernel) multiplicative VCG identifies the correct
mistake step with probability tending to one. We sketch the analogue for
the general power kernel.

\begin{theorem}[Power-VCG Anchored Correctness, sketch]
\label{thm:power-anchored}
Suppose the discriminator reports are $(f, \gamma, \eta)$-anchored with
$f \leq \lfloor (K{-}1)/2 \rfloor$ and $\eta < c$, and the margin
conditions $m_{\mathrm{thr}}, m_{\mathrm{step}}, m_{\mathrm{agent}} > 0$
of the theory note hold. Let $\db^*$ be the Power-VCG allocation
(Algorithm~\ref{alg:power}). Then there exist constants $C_1(p, c), c'(p, c) > 0$
depending only on $(K, \varepsilon, p, c)$ such that for every $t \in [T]$,
\[
|d_t^* - \theta_t| \;\leq\; \eta + C_1(p, c) \exp\!\big(-c'(p, c)\,\gamma\,(T-1)\big),
\]
and the attribution decisions are correct whenever the right-hand side
is less than the corresponding margin.
\end{theorem}

\paragraph{Proof idea.}
The proof is structurally the same as for $p = 2$, with two changes:
(i)~the active-set lower bound on anchor weights becomes
$(1 - \eta^p)^{(1-\gamma)(T-1)}(1 - (\eta + c)^p)^{\gamma(T-1)}$ rather
than the squared-loss expression, decoupling robustness from report scale
whenever $p < 2$;
(ii)~the FOC robustness analysis uses the M-estimator property:
the $p$-quantile is robust to any coalition of weight $< 1/2$, which holds
for all $p \geq 1$ but with different breakdown constants. The exponent
$p$ enters the bound through $C_1, c'$ in a smooth way, so the theorem
applies uniformly across the family. The corresponding full proof,
generalising Tong's L1-Huber-specific argument in the previous note
revision, is left for the companion theory-note revision.

\begin{remark}[Implications for the operating points]
At $p = 1$, the constants $C_1, c'$ depend only on $(c, \eta, \gamma)$, not
on the report scale, which explains the strong HC performance.
At $p = 1.75$, $c = 1$ (the AG operating point, no truncation), the
constants depend on report scale but the smoother kernel produces
better-conditioned coordinate updates when reports are bounded away from
saturation. Theorem~\ref{thm:power-anchored} subsumes both regimes within
one theoretical framework.
\end{remark}

\section{Limitations and Open Questions}
\label{sec:limitations}

\paragraph{The two-element menu is data-driven.}
The two operating points $(p_L, c_L, \varepsilon_L)$ and
$(p_S, c_S, \varepsilon_S)$ in Algorithm~\ref{alg:adaptive} are chosen
from the empirical $(p, c)$ scan on Who\&When. While the kernel family is
principled, the specific menu is dataset-specific. A theory that
\emph{predicts} the optimal $(p, c)$ from observable noise structure
(e.g.\ from the empirical distribution of $\thetahat$) is open.

\paragraph{Truth-aligned OMD calibration.}
Section~\ref{sec:omd-diagnostic} reports that three OMD objectives
(payment, entropy, margin) all degrade AG step accuracy by a common
mechanism: consensus drift moves the predicted step toward the
highest-agreement location, which on AG diverges from the off-consensus
mistake step. A principled follow-up is OMD on a non-consensus objective.
One concrete candidate: identify, in each trace, the discriminator most
often \emph{opposed} to the running allocation (using cross-step or
cross-trace statistics), and gradient-ascent on agreement with that
discriminator's argmax. This is a heuristic but breaks the symmetry that
makes the existing objectives biased. A formal characterisation of
"truth-aligned" calibration signals is an open problem; we conjecture
that they cannot be defined purely in terms of the within-trace
welfare~$V^{(p,c)}$, because all welfare-derived quantities inherit the
consensus bias documented in Section~\ref{sec:omd-diagnostic}.

\paragraph{Trace length as input is benchmark-flexible.}
The fact that $T$ separates HC from AG with $3.6\%$ error is a property
of Who\&When's two-subset structure (AG was constructed algorithmically
with $T = 10$). The smooth selector~\eqref{eq:smooth-selector} accommodates
benchmarks with continuous trace-length distributions: setting the
smoothing scale $\tau$ to a moderate value (e.g.\ $0.3$--$1.0$) produces
soft interpolation between the two anchor configurations on traces of
intermediate length. On benchmarks where trace length alone is insufficient
to distinguish regimes (e.g.\ when all traces have similar lengths but
heterogeneous noise structure), the combined-score extension
$s(\thetahat) = \log T + \alpha \cdot \mathrm{kurt}(\thetahat)$ from
Section~\ref{sec:combined} provides a natural generalisation. On future
benchmarks, $(\mu, \tau, \alpha)$ can be tuned via held-out validation;
the underlying kernel-family construction $\kappa_{c, p}$ and the
M-estimator coordinate FOC are benchmark-independent.

\paragraph{Larger $K$.}
With $K = 3$ and $f = 2$, the anchoring hypothesis of
Theorem~\ref{thm:power-anchored} fails on AG (only one discriminator
is genuinely informative). With $K = 5$ or $K = 7$, anchoring becomes
more plausible and the dominance regime may dissolve. Re-running with
two additional discriminators is the most direct next test.

\section{Discussion Points for Collaborator Review}

\begin{enumerate}[leftmargin=*]
\item Should we extend the experiments to $K = 5$ before submission, or
treat that as follow-up?
\item Theorem~\ref{thm:power-anchored} full proof: main paper, appendix,
or a revision of the companion theory note?
\item Proposition~\ref{prop:gs-nash} (Gauss--Seidel coordinate-wise Nash):
should this be elevated to a theorem with the additional uniqueness
argument under the strict monotonicity hypothesis for $p \in (1, 2]$, or
kept as a proposition with the proof we have?
\item Should we attempt to derive the optimal $(p, c)$ predictively from
report-distribution statistics, or report the data-driven choice as is?
\item OMD diagnostic (Section~\ref{sec:omd-diagnostic}): include as a
section in the main paper to motivate $R = 0$, or move to an appendix?
\item The truth-aligned OMD direction suggested in
Section~\ref{sec:limitations}: pursue as a follow-up paper, or attempt a
preliminary experiment for this submission?
\item The mechanism interpretation: it is convenient to describe
Algorithm~\ref{alg:adaptive} as "one mechanism with a two-element
configuration menu" rather than "two mechanisms with a router". Is this
framing precise enough, or should we make the menu structure more
explicit in the paper?
\end{enumerate}

\appendix

\section{Reproducibility}

All numbers reproduce with the scripts in the project repository.
\begin{sloppypar}
\begin{itemize}[nosep]
\item Per-method accuracy on the squared-kernel baseline (Section~3.1):
\texttt{scripts/run\_vcg.py} verified to give
$\mathrm{Acc}_{\mathrm{step}} = 0.349$ on AG with $R = 5$,
$\varepsilon_{\text{report}} = 0.10$.
\item $(p, c)$ scan (Section~3.2): \texttt{scripts/scan\_pc.py}.
\item Within-trace regime statistics (Section~3.4):
\texttt{scripts/regime\_stats.py}.
\item Smooth trace-length-adaptive selector (Section~3.3):
\texttt{scripts/adaptive\_kernel\_smooth.py}
(default mode \texttt{T-only} with
$\mu = \log 10$, $\tau = 0.01$ recovers the canonical operating point;
\texttt{--scan} sweeps $(\mu, \tau)$ for sensitivity analysis).
The hard-threshold variant in
\texttt{scripts/adaptive\_kernel.py}
is retained for the $\tau \to 0^+$ limit and yields identical accuracy.
\item Combined-score selector with kurtosis (Section~3.5):
\texttt{scripts/adaptive\_kernel\_combined.py},
where \texttt{--scan} sweeps $(\alpha, \mu, \tau)$.
\item Power-kernel allocation: \texttt{vcg/allocation\_power.py}.
\item OMD calibration with selectable objective
(Section~\ref{sec:omd-diagnostic}):
\texttt{scripts/run\_power\_v2.py --omd-objective \{payment, entropy, margin\}}.
The corresponding pipeline modules \texttt{vcg/pipeline\_power.py},
\texttt{vcg/gradient\_power.py}, \texttt{vcg/payment\_power.py},
and \texttt{vcg/gradient\_concentration.py} implement the three
gradient signals.
\item Per-trace OMD diagnostic (which traces had argmax changed by OMD,
hit/miss flips): \texttt{scripts/omd\_find\_changers.py}.
\end{itemize}
\end{sloppypar}
Random seed $42$ throughout; the $(p, c)$ scan and adaptive selector are
deterministic (no random component).

\end{document}
